\chapter[Requerimientos de red - Software]{Requerimientos de red - Software}
\label{cp:requerimientos-software}

\section*{Resumen}
Los laboratorios de enseñanza del Departamento de Física cuentan con 5 espacios/laboratorios con entre 12 y 16 computadoras cada uno, las cuales se necesita que estén conectadas a una red local para poder compartir recursos y acceder a internet. Todas estas computadoras, al día de hoy, usan Windows como sistema operativo. Se intentó hacer un cambio a Linux, pero no se pudo por falta de soporte para muchos de los equipos y la falta de conocimiento en el manejo de este por parte del cuerpo docente y el personal técnico de los laboratorios.

Muchos de los laboratorios, además de necesitar bocas de red para las computadoras, necesitan bocas de red para otros equipos como proyectores, impresoras, equipos de medición, etc.

En esta parte del artículo solo se hará referencia a los requerimientos de software. 

\section{Necesidades de recursos compartidos y red}
\subsection{Permisos}
Cuando llegué al laboratorio (y desde que usé las PC como alumno, si no recuerdo mal), los alumnos tienen un usuario general de acceso a las PC con el que pueden usar con libertad el software y guardar archivos en las carpetas de alumnos, Documentos y Descargas, aunque también tienen el mal hábito de usar el escritorio para esto. Por otro lado, no tienen ningún privilegio para instalar o borrar programas, ni para cambiar cualquier configuración. 
Hasta donde me comentaron, esto es porque de otra manera el mantenimiento era inviable o al menos esto es lo que siempre se argumentó.

Por otro lado, el personal de los laboratorios, que anteriormente hacía mantenimiento básico y hoy genera las imágenes e instala/reinstala todo el software, debe contar con muchos más privilegios. Sobre todo, privilegios de instalación de software:

\begin{description}
    \item[Alumnos:] Usuario y contraseñas únicas. Uso de software y guardado de archivos en lugares seleccionados.
    \item[Docentes:] Usuario y contraseñas únicas. En principio igual que los alumnos.
    \item[Técnico:] Usuario y contraseñas por usuario. Acceso de administrador.
\end{description}

Existen algunos problemas con el acceso de escritura restringido para alumnos. Algunos programas como Anaconda para crear los entornos virtuales, por defecto, usan un subdirectorio en su carpeta de instalación, y al estar restringido el acceso, la creación de entornos virtuales (algo muy útil para alumnos) falla.

\subsection{Impresoras} 
Todos los laboratorios cuentan con impresoras locales pero, por falta de hojas o mal funcionamiento, es habitual que los alumnos necesiten acceder a las impresoras de otro laboratorio. Por esto es necesario que las computadoras de todos los laboratorios puedan acceder a las impresoras de todos los laboratorios.

\subsection{Archivos} 
Algo similar pasa entre los archivos de las computadoras. Tanto para instalar imágenes o programas pesados, como para compartir archivos entre los alumnos, es necesario que las computadoras de todos los laboratorios puedan acceder a los archivos de todos los laboratorios. Antiguamente se habían creado carpetas de uso general (una para cada laboratorio) para que los alumnos puedan guardar y acceder a los archivos en ese lugar; si no recuerdo mal, todas las computadoras de cada laboratorio montaban una unidad de almacenamiento de red en el inicio de sesión para que la carpeta compartida sea fácil de acceder. En otro momento también se compartió la carpeta de Documentos de cada computadora de alumnos como un recurso de red, y eso facilitó mucho el acceso a los archivos. No tengo idea de cuál de las soluciones es mejor. 
